\exercisegroup

\begin{enumerate}
\def\labelenumi{\arabic{enumi})}
\tightlist
\item
  设 \((\mathrm{C}_{i})_{i\in\mathrm{I}}\) 是一族范畴。设 C 是诸范畴 \(C_{i}\) 的底箭图组成的族的和箭图。
\end{enumerate}

\begin{enumerate}
\def\labelenumi{\alph{enumi})}
\item
  证明箭图 C 赋以由诸范畴 \(C_{i}\) 的合成法则诱导的合成法则后是一个范畴。我们称之为族 \(C_{i}\) 的和范畴。
\item
  证明从 \(C_{i}\) 到 C 的典范态射是一个函子。
\item
  若诸 \(C_{i}\) 都是群胚，证明范畴 C 是一个群胚。
\end{enumerate}

\begin{enumerate}
\def\labelenumi{\arabic{enumi})}
\setcounter{enumi}{1}
\tightlist
\item
  设 \((\mathrm{C}_{i})_{i\in\mathrm{I}}\) 是一族范畴。设 C 是诸范畴 \(C_{i}\) 的底箭图的积。
\end{enumerate}

\begin{enumerate}
\def\labelenumi{\alph{enumi})}
\item
  证明箭图 C 赋以由诸范畴 \(C_{i}\) 的合成法则诱导的合成法则后是一个范畴。称之为族 \((\mathrm{C}_{i})\) 的积范畴。
\item
  证明典范态射 \(pr_{i}: C \rightarrow C_{i}\) 是一个函子。
\item
  若诸 \(C_{i}\) 都是群胚，证明范畴 C 是一个群胚。
\end{enumerate}

\begin{enumerate}
\def\labelenumi{\arabic{enumi})}
\setcounter{enumi}{2}
\tightlist
\item
  设 C 是一个范畴；用 m 表示它的合成法则。设 C° 是反有向图，赋以由 \(m^{\circ}(f,g)=m(g,f)\)（对 C 的每对可合成的箭 \((f,g)\)）给出的合成法则。
\end{enumerate}

\begin{enumerate}
\def\labelenumi{\alph{enumi})}
\item
  证明 C° 是一个范畴。称之为 C 的反范畴。
\item
  若 C 是一个群胚，证明 C° 是一个同构于 C 的群胚。
\end{enumerate}

\begin{enumerate}
\def\labelenumi{\arabic{enumi})}
\setcounter{enumi}{3}
\tightlist
\item
  我们把群胚 G 的质量 \(\mu(\mathrm{G})\) 定义为诸和 \(\sum_{a\in\mathrm{A}}\mathrm{Card}(\mathrm{G}_{a})^{-1}\) 的上确界，其中 A 取遍 \(\mathrm{Som}(\mathrm{G})\) 中与 G 的每个轨道至多交于一点的子集所成的集合。若 \(\mu(\mathrm{G})\in\mathbf{R}\)，称 G 具有有限质量；若 G 的轨道的集合有限且它在每点的迷向群有限，称 G 是本质有限的。
\end{enumerate}

\begin{enumerate}
\def\labelenumi{\alph{enumi})}
\item
  证明一个本质有限的群胚具有有限质量。
\item
  设 G 是与有限群 \(\Gamma\) 相伴的群胚。它是本质有限的，且有 \(\mu(\mathrm{G}) = 1/\mathrm{Card}(\Gamma)\)。
\item
  设 \(\Gamma\) 是一个有限群，X 是一个赋有 \(\Gamma\) 的右作用的有限集，G 是相伴的群胚。它是本质有限的，且有 \(\mu(\mathrm{G}) = \mathrm{Card}(\mathrm{X}) / \mathrm{Card}(\Gamma)\)。
\end{enumerate}

\(d)\) 设 \((\mathrm{G}_i)_{i\in \mathrm{I}}\) 是一族具有有限质量的群胚。族 \((\mathrm{G}_i)\) 的和群胚与积群胚的质量分别为 \(\sum \mu (\mathrm{G}_i)\) 与 \(\prod \mu (\mathrm{G}_i)\)。

\begin{enumerate}
\def\labelenumi{\alph{enumi})}
\setcounter{enumi}{4}
\tightlist
\item
  设 X 是一个集合。我们定义一个群胚 G，其顶点集是 X 的有限子集组成的集合，且对 X 的每对有限子集 \((Y_{1}, Y_{2})\)，以 \(Y_{1}\) 为起点、\(Y_{2}\) 为终点的箭的集合是从 \(Y_{1}\) 到 \(Y_{2}\) 的双射组成的集合，箭的合成由映射的合成给出。证明 G 是本质有限的，并计算它的质量。
\end{enumerate}

\begin{enumerate}
\def\labelenumi{\arabic{enumi})}
\setcounter{enumi}{4}
\tightlist
\item
  \begin{enumerate}
  \def\labelenumii{\alph{enumii})}
  \tightlist
  \item
    设 \(\varphi\colon G\to G'\) 是一个群胚态射，使得映射 \(\operatorname{Som}(\varphi)\) 是单射。证明 \(\varphi(G)\) 是 \(G'\) 的一个子群胚。
  \end{enumerate}
\end{enumerate}

\begin{enumerate}
\def\labelenumi{\alph{enumi})}
\setcounter{enumi}{1}
\tightlist
\item
  设 G 是一个非传递群胚，其顶点集 \(\operatorname{Som}(G)\) 是二元集 \(\{a,b\}\)。设 G' 是与自由积 \(G_a*G_b\) 相伴的群胚。存在唯一的群胚态射 \(\varphi:G\to G'\)，使得映射 \(\varphi_a\) 与 \(\varphi_b\) 分别是从 \(G_a\) 与 \(G_b\) 到 \(G_a*G_b\) 的典范映射。若 \(G_a\) 或 \(G_b\) 不是平凡群，则该群胚态射的像 \(\varphi(G)\) 不是 \(G'\) 的子群胚。
\end{enumerate}

\begin{enumerate}
\def\labelenumi{\arabic{enumi})}
\setcounter{enumi}{5}
\tightlist
\item
  设 G 与 G' 是图，\(\varphi\colon\mathrm{G}'\to\mathrm{G}\) 是一个图态射。我们称 \((\mathrm{G}',\varphi)\) 是 G 的一个覆叠，如果对每个顶点 \(a\in\mathrm{Som}(\mathrm{G}')\) 及每条箭 \(f\in\mathrm{Fl}(\mathrm{G})\)（以 \(\mathrm{Som}(\varphi)(a)\) 为起点），存在唯一的箭 \(f'\in\mathrm{Fl}(\mathrm{G}')\)（以 \(a\) 为起点）使 \(\mathrm{Fl}(\varphi)(f')=f\)。
\end{enumerate}

我们假设 \((\mathrm{G}^{\prime},\varphi)\) 是 G 的一个覆叠。

\begin{enumerate}
\def\labelenumi{\alph{enumi})}
\item
  证明存在 \(\operatorname{Grp}(G)\) 在 \(\operatorname{Som}(G')\) 上的、相对于映射 \(\operatorname{Som}(\varphi)\) 的唯一作用，使得 \(x \cdot \operatorname{Fl}(\varphi)(f) = t(f)\)（对每个顶点 \(x \in \operatorname{Som}(G')\) 及每条以 x 为起点的箭 \(f \in \operatorname{Fl}(G')\)）。
\item
  由此推出，若 x 与 y 是 G 的属于同一连通分支的顶点，则 \(\text{Card}(\text{Som}(\varphi)^{-1}(x)) = \text{Card}(\text{Som}(\varphi)^{-1}(y))\)。
\end{enumerate}

\begin{enumerate}
\def\labelenumi{\arabic{enumi})}
\setcounter{enumi}{6}
\tightlist
\item
  设 G 是一个图，\((G_{1},\varphi_{1})\) 与 \((G_{2},\varphi_{2})\) 是图 G 的覆叠。
\end{enumerate}

\begin{enumerate}
\def\labelenumi{\alph{enumi})}
\item
  设 \(\psi\colon\mathrm{G}_{1}\to\mathrm{G}_{2}\) 是一个满足 \(\varphi_{2}\circ\psi=\varphi_{1}\) 的图态射。证明有 \(\psi(x\cdot\gamma)=\psi(x)\cdot\gamma\)（对所有 \(x\in\mathrm{Som}(\mathrm{G}_{1})\) 及每条箭 \(\gamma\)（属于 \(\varpi_{\mathrm{G}}\)，起点为 \(\varphi_{1}(x)\)））。
\item
  设 \(u \colon \operatorname{Som}(\mathrm{G}_1) \to \operatorname{Som}(\mathrm{G}_2)\) 是一个映射，满足 \(u(x \cdot \gamma) = u(x) \cdot \gamma\)（对所有 \(x \in \operatorname{Som}(\mathrm{G}_1)\) 及每条箭 \(\gamma\)（属于 \(\varpi_{\mathrm{G}}\)，起点为 \(\varphi_1(x)\)））。证明存在唯一的图态射 \(\psi \colon \mathrm{G}_1 \to \mathrm{G}_2\)，使得 \(\varphi_2 \circ \psi = \varphi_1\) 且 \(\operatorname{Som}(\psi) = u\)。
\item
  设 X 是一个集合，\(p: X \to \text{Som}(G)\) 是一个映射。假设给定了 \(\varpi_{G}\) 相对于 p 在 X 上的一个作用。证明存在一个以 X 为顶点集的图 \(G'\) 及一个图态射 \(\varphi: G' \to G\)，使得 \(\text{Som}(\varphi) = p\)，\((G', \varphi)\) 是 G 的一个覆叠，且 \(\varpi_{G}\) 在 \(\text{Som}(G')\) 上的作用与给定的作用一致。
\end{enumerate}

\begin{enumerate}
\def\labelenumi{\arabic{enumi})}
\setcounter{enumi}{7}
\tightlist
\item
  设 G 是一个连通图，\(a\) 是 G 的一个顶点。
\end{enumerate}

\begin{enumerate}
\def\labelenumi{\alph{enumi})}
\tightlist
\item
  设 \((\mathrm{G}_{1},\varphi_{1})\) 与 \((\mathrm{G}_{2},\varphi_{2})\) 是图 G 的覆叠。设 \(u\colon\operatorname{Som}(\mathrm{G}_{1})\to\operatorname{Som}(\mathrm{G}_{2})\) 是一个满足 \(u(x)\cdot\gamma=u(x\cdot\gamma)\) 的映射
\end{enumerate}

（对所有 \(\gamma \in \varpi_{\mathrm{G},a}\) 及每个顶点 \(x\)（属于 \(\mathrm{G}_1\) 且满足 \(\operatorname{Som}(\varphi_1)(x) = a\)））。证明存在唯一的图态射 \(\psi: \mathrm{G}_1 \to \mathrm{G}_2\) 使得 \(\psi = \operatorname{Som}(u)\)。

\begin{enumerate}
\def\labelenumi{\alph{enumi})}
\setcounter{enumi}{1}
\tightlist
\item
  设 X 是一个赋有群 \(\varpi_{\mathrm{G},a}\) 的作用的集合。证明存在图 G 的一个覆叠 \((\mathrm{G}',\varphi)\)，使得 \(\operatorname{Som}(\varphi)^{-1}(a) = \mathrm{X}\)，且 \(\varpi_{\mathrm{G}}\) 在 \(\operatorname{Som}(\mathrm{G}')\) 上的作用诱导出 \(\varpi_{\mathrm{G},a}\) 在 X 上给定的作用。
\end{enumerate}

\begin{enumerate}
\def\labelenumi{\arabic{enumi})}
\setcounter{enumi}{8}
\tightlist
\item
  设 G 与 G' 是图，\(\varphi\colon G'\to G\) 是一个图态射，使得 \((G',\varphi)\) 是 G 的一个覆叠。
\end{enumerate}

\begin{enumerate}
\def\labelenumi{\alph{enumi})}
\item
  证明图 \(\varphi(\mathrm{Grp}(\mathrm{G}^{\prime}))\) 是群胚 \(\mathrm{Grp}(\mathrm{G})\) 的一个子群胚。
\item
  证明对每个顶点 \(a\)（属于 \(\mathbf{G}'\)），群同态 \(\pi_1(\varphi, a)\) 是单射。
\item
  设 G 连通，并固定 G 的一个顶点 a。设 K 是 \(\varpi_{G,a}\) 的一个子群。证明存在图 G 的一个覆叠 \((H,\psi)\) 及 H 的一个顶点 b，使得 \(\psi(b)=a\) 且 K 是态射 \(\pi_{1}(\psi,a)\) 的像。
\item
  证明自由群的子群是自由的（另见 IV, p. 417, 例 2）。
\item
  设 n 与 q 是整数；若 F 是 n 个生成元上的自由群，而 K 是 F 的指数为 q 的子群，证明 K 是 \(1 + q(n - 1)\) 个生成元上的自由群。
\end{enumerate}

\begin{enumerate}
\def\labelenumi{\arabic{enumi})}
\setcounter{enumi}{9}
\tightlist
\item
  设 F 是一个集合，C 是 \(F \times F\) 的一个子集，\(m: C \to F\) 是一个映射。我们作如下假设：
\end{enumerate}

\begin{enumerate}
\def\labelenumi{(\roman{enumi})}
\item
  由 \((f,g)\mapsto(f,m(f,g))\) 与 \((f,g)\mapsto(m(f,g),g)\) 给出的从 C 到 F × F 的映射都是单射；
\item
  若 \(f, g, h \in \mathrm{F}\) 使得对 \((f, g)\) 与 \((g, h)\) 属于 C，则对 \((f, m(g, h))\) 与 \((m(f, g), h)\) 属于 C，且有 \(m(f, m(g, h)) = m(m(f, g), h)\)。
\item
  对每个 \(f \in \mathrm{F}\)，存在元素 \(o(f)\)、\(t(f)\) 与 \(\overline{f} \in \mathrm{F}\)（由 (i) 知必然唯一），使得对 \((f, t(f))\)、\((o(f), f)\)、\((f, \overline{f})\) 与 \((\overline{f}, f)\) 属于 C，且有 \(m(f, t(f)) = f = m(o(f), f)\) 及 \(m(f, \overline{f}) = o(f)\)。
\end{enumerate}

\begin{enumerate}
\def\labelenumi{\alph{enumi})}
\item
  证明对每个 \(f \in F\)，有 \((\overline{f}, f) \in C\) 且 \(m(\overline{f}, f) = t(f)\)。由此推出 \((o(f), o(f)) \in C\) 且 \(m(o(f), o(f)) = o(f)\)。再证明 \(o(\overline{f}) = t(f)\) 且 \(\overline{\overline{f}} = f\)。
\item
  设 S 是 F 中元素 \(e\)（满足 \((e,e)\in\mathrm{C}\) 且 \(m(e,e)=e\)）组成的集合；用 \(o\) 与 \(t\) 表示由 \(f\mapsto o(f)\) 与 \(f\mapsto t(f)\) 给出的从 F 到 S 的映射。
\end{enumerate}

证明四元组 \(\Gamma = (S, F, o, t)\) 是一个箭图，且映射 \(m\) 是 \(\Gamma\) 上的一个使它成为群胚的合成法则。（这一构造属于 H. BRANDT；参见 “Über eine Verallgemeinerung des Gruppenbegriffes”, Mathematische Annalen (1926), vol. 96, p. 360–366。）

\begin{enumerate}
\def\labelenumi{\arabic{enumi})}
\setcounter{enumi}{10}
\tightlist
\item
  设 G 是一个群胚；设 C 是 G 的可合成的箭对组成的集合，\(m\colon\mathrm{C}\to\mathrm{Fl}(\mathrm{G})\) 是 G 的合成法则。
\end{enumerate}

\begin{enumerate}
\def\labelenumi{\alph{enumi})}
\item
  证明 C 与 m 满足习题 10 的假设 (i)、(ii)、(iii)。用 Γ 表示该习题中的构造所给出的群胚。
\item
  证明由映射 \(e \mapsto o(e)\)（从 \(\operatorname{Som}(\Gamma)\) 到 \(\operatorname{Som}(\mathrm{G})\)）与集合 \(\operatorname{Fl}(\mathrm{G})\) 到其自身的恒同映射组成的对是一个从 \(\Gamma\) 到 G 的群胚同构。
\end{enumerate}

这表明一个群胚由它的箭集及其合成法则确定。

\begin{enumerate}
\def\labelenumi{\arabic{enumi})}
\setcounter{enumi}{11}
\tightlist
\item
  设 G 是一个群，S 是 G 的子群组成的集合，F 是 G 中满足存在 G 的子群 H 与 G 的元素 g 使 X = gH 的子集 X 组成的集合。
\end{enumerate}

\begin{enumerate}
\def\labelenumi{\alph{enumi})}
\item
  设 \(\mu\) 是从 \(G \times G \times G\) 到 \(G\) 的、由 \(\mu(x,y,z)=xy^{-1}z\) 定义的映射。证明一个非空子集 \(X\)（\(G\) 的）属于集合 F 当且仅当 \(\mu(X\times X\times X)\subset X\)。
\item
  证明取 X 的任意元素 g 并令 \(o(\mathrm{X}) = g^{-1}\mathrm{X}\) 与 \(t(\mathrm{X}) = \mathrm{X}g^{-1}\)，便定义了映射 \(o: F \to S\) 与 \(t: F \to S\)。
\item
  若 \(X_{1}\) 与 \(X_{2}\) 是 S 中满足 \(t(\mathrm{X}_{1}) = o(\mathrm{X}_{2})\) 的元素，令 \(m(\mathrm{X}_{1}, \mathrm{X}_{2}) = \mathrm{X}_{1}\mathrm{X}_{2}\)。证明映射 m 是箭图 (S,F,o,t) 上的一个使它成为群胚的合成法则（“Baer 群胚”）。把它记作 \(\mathrm{Ba}(\mathrm{G})\)。
\item
  证明 \(\mathrm{Ba}(\mathrm{G})\) 的轨道是 G 的子群的共轭类。
\item
  证明 Ba(G) 在 S 的元素 H 处的迷向群同构于商群 N(H)/H。
\end{enumerate}

\setcounter{exercisegroup}{4}

